\documentclass{../../../unipd-notes}

\unipdsetup{
  course = {Fondamenti di Ingegneria Informatica},
  short-course = {Fondamenti di Ingegneria},
  professor = {Prof. Marco Rossi},
  academic-year = {2026--2027},
  degree-year = {1},
  semester = {1},
  document-type = {Appunti delle lezioni},
  date = {3 agosto 2026},
  version = {1.0.0}
}

\begin{document}
\chapter*{Informazioni sul corso}

La seguente scheda identifica l'insegnamento, l'edizione degli appunti e il contesto accademico di riferimento.

\begin{tabularx}{\textwidth}{>{\bfseries\sffamily}lL}
  Corso & \unipdmetadata{course} \\
  Nome breve & \unipdmetadata{short-course} \\
  Docente & \unipdmetadata{professor} \\
  Anno accademico & \unipdmetadata{academic-year} \\
  Anno di corso & \unipdmetadata{degree-year} \\
  Semestre & \unipdmetadata{semester} \\
  Tipo di documento & \unipdmetadata{document-type} \\
  Autore & \unipdmetadata{author} \\
  Università & \unipdmetadata{university} \\
  Corso di laurea & \unipdmetadata{degree} \\
  Repository & \unipdmetadata{repository} \\
  Licenza & \unipdmetadata{license} \\
\end{tabularx}

\begin{remark}
Gli appunti seguono l'ordine delle lezioni; eventuali correzioni sono riportate nella cronologia delle revisioni.
\end{remark}
\end{document}
